\documentclass[10pt,a4paper]{article}
\usepackage[margin=15mm,top=14mm,bottom=16mm]{geometry}
\usepackage[T1]{fontenc}
\usepackage[utf8]{inputenc}
\usepackage{lmodern,microtype,amsmath,amssymb,graphicx,booktabs,array,xcolor}
\usepackage[font=small,labelfont=bf,skip=4pt]{caption}
\usepackage[colorlinks=true,urlcolor=blue!50!black,linkcolor=blue!50!black]{hyperref}
\hypersetup{pdftitle={Singlet coverings in two and three dimensions},
  pdfsubject={Square--octagon RVB, octochlore coverings, and a possible layer construction}}
\usepackage{fancyhdr}
\pagestyle{fancy}
\fancyhf{}
\renewcommand{\headrulewidth}{0pt}
\fancyfoot[L]{\footnotesize Singlet coverings: collaborator notes \quad 22 October 2026}
\fancyfoot[R]{\footnotesize\thepage/4}
\setlength{\parindent}{0pt}
\setlength{\parskip}{5pt}
\setlength{\abovedisplayskip}{5pt}
\setlength{\belowdisplayskip}{5pt}
\setlength{\textfloatsep}{5pt}
\newcommand{\sect}[1]{\vspace{5pt}{\large\bfseries\color{blue!45!black}#1}\par\vspace{2pt}}
\newcommand{\ket}[1]{\lvert#1\rangle}
\newcommand{\eps}{\epsilon}
\newcommand{\xor}{\mathbin{\oplus}}
\begin{document}

{\LARGE\bfseries Singlet coverings in two and three dimensions}\par
{\normalsize Square--octagon RVB, octochlore coverings, and a possible layer construction}

\sect{1. Two-dimensional model and the two dummy-link arrangements}
The underlying graph is the trivalent square--octagon lattice. Its gray squares $P$ are
drawn upright; a physical spin $1/2$ sits at the midpoint of each square edge
(blue). Each dummy link carries \emph{two} distinct endpoint spins (orange).
The three slots around a graph vertex form a triangle, which contains zero or
one singlet. Every physical spin is covered exactly once globally.

\begin{center}
\includegraphics[width=.87\linewidth]{upright-geometry.pdf}
\end{center}
\vspace{-6pt}
\captionof{figure}{Gray lines define the base graph; green triangle edges are allowed
singlet bonds. The gray squares have blue spins on their edges. A cell contains
four blue spins and four orange endpoint spins.}

On a dummy link either its two endpoint spins form an \textbf{internal singlet},
or \textbf{both endpoints pair into their own triangles}. No spin is left free.
In the internal case each triangle is empty or contains its blue--blue singlet:
four possibilities. In the external case each endpoint chooses one of two blue spins:
another four.

\begin{center}
\includegraphics[width=\linewidth]{dummy-covering-table.pdf}
\end{center}
\vspace{-6pt}
\captionof{table}{Eight local coverings and square-vertex loop patterns. Triangles and
gray square-edge corners follow the lattice orientation. Solid green lines are singlets;
dashed lines are triangle frames; dotted stubs pair blue spins in other triangles.
After contraction $a=N$, $b=W$, $c=E$, $d=S$. Bits refer to the local reference
$0000$ with the orange spins internally paired; this extends to a global covering.}

Writing $a,b,c,d=1$ when an ordinary slot is used in its displayed triangle, and
$t=1$ when a dummy endpoint is used there, the local constraints are
\[
 a\xor b=t=c\xor d,
 \qquad a\xor b\xor c\xor d=0.
\]
In this representative reference these are also the loop bits. For any other
reference, first take the difference of ownership bits modulo two.

The conventional RVB state sums the normalized singlet products with equal coefficients:
\[
 \ket{\Psi_{2D}}=\mathcal N^{-1/2}\sum_{D\ {
m allowed}}\ket{D},
 \qquad
 \ket{s_{ij}}=\frac{\ket{\uparrow_i\downarrow_j}-\ket{\downarrow_i\uparrow_j}}{\sqrt2}.
\]
Singlet arrows run counterclockwise around each triangle; dummy arrows follow
the positive simulation axes. This fixes relative phases. No additional loop-number
sign is included. Equal coefficients do not make overlapping singlet coverings orthogonal.

\newpage
\sect{2. Coverings $\longrightarrow$ loops $\longrightarrow$ the square lattice}
Choose one valid reference covering. An ordinary-edge spin has exactly one owner;
the two dummy endpoints are either both used by triangles or neither is. Ownership
differences agree at both ends of every edge and have even parity at each base vertex.
Since the base graph has degree three, the selected edges form closed loops [1].

\begin{center}
\includegraphics[width=\linewidth]{corrected-mapping-upright.pdf}
\end{center}
\vspace{-6pt}
\captionof{figure}{One checked covering and its reference-relative loops. The relevant
gray squares are upright in \emph{all three panels}. The reference has internal dummy
singlets and alternating occupied/empty triangles around each square. Boundary singlet
stubs continue periodically. Dashed green triangles mark the allowed singlet supports.
Dummy edges shrink to vertices; square edges survive.}

Here ``cutting/erasing'' a dummy edge means \textbf{contracting its two graph endpoints},
not deleting its physical spins from the RVB state. The result is the regular square
graph. Writing $n$ for reference-difference bits, the erased bit is recovered as
$n_{\rm dummy}=n_a\xor n_b=n_c\xor n_d$, so the map is a
bijection of even-edge configurations, including winding sectors. At a four-leg
vertex retain the inherited $a$--$b$, $c$--$d$ routing shown in Table~1.
These reference loops label coverings; bra--ket transition loops instead determine
singlet overlaps. Consequently the map alone does not determine quantum correlations.

\sect{3. Bond correlations: numerical evidence and diagnosis}
We measure the connected physical bond covariance
$C_B(r)=\langle B_0B_r\rangle-\langle B_0\rangle\langle B_r\rangle$,
$B=\mathbf S_i\!\cdot\!\mathbf S_j$, for a square-edge-spin bond inside a triangle.
It equals the connected covariance of the singlet projector $1/4-B$.

\begin{center}
\includegraphics[width=\linewidth]{correlations-2d.pdf}
\end{center}
\vspace{-6pt}
\captionof{figure}{Our model's computed bond correlations, not published curves.
An exact $D=3$ PEPS is contracted with infinite-plane CTMRG ($\chi=64$), comparing
two truncation cutoffs. Shading at left marks the fit window $r=12$--18.}

Following PEPS boundary and sector diagnostics [2,3], we compare cutoffs, fit windows,
and exact finite-cylinder checks. Within each fixed parity sector, transfer modes
coupled to $B$ give $\xi_a=-1/\log|\lambda_a/\lambda_0|$.
Fits to $Ae^{-r/\xi}$ and $Ar^{-p}e^{-r/\xi}$ give
$\xi=0.82$ and $0.86$ cells. The local inverse length
$\kappa_{\rm eff}=\log|C_B(r)/C_B(r+1)|$ tends to $1/\xi>0$ for an exponential
and to zero for a pure power law. The resolved data support a short exponential-like
scale, \textbf{$\xi\sim0.8$--$0.9$ simulation-cell translations}; an 8\% cutoff
change at $r=18$ prevents a precise bulk claim. One translation is a diagonal
$(1,-1)$ step on the upright contracted grid. Orthogonal-covering occupation
correlations instead have exactly finite range; they are a different observable and state.

{\footnotesize
[1] M.~Iqbal, D.~Poilblanc and N.~Schuch, \emph{Semionic resonating valence bond states}
(2014), \href{https://arxiv.org/abs/1407.7773}{arXiv:1407.7773}. Reference-arrow construction;
our state has no semionic sign.\par
[2] N.~Schuch et al., \emph{Resonating valence bond states in the PEPS formalism}
(2012), \href{https://arxiv.org/abs/1203.4816}{arXiv:1203.4816}.\par
[3] D.~Poilblanc et al., \emph{Topological and entanglement properties of resonating
valence bond wave functions} (2012), \href{https://arxiv.org/abs/1202.0947}{arXiv:1202.0947}.
The papers motivate the diagnostics; they do not calculate this square--octagon model.}

\newpage
\sect{4. Three-dimensional octochlore covering model}
This section incorporates the geometric model and cage-net correspondence of
J.~Garre Rubio's supplied notes (19 October 2026), so they need not be read separately.
The corner-sharing octahedra have simple-cubic centres $v\in\mathbb Z^3$ [4]:
\[
 O_v=\operatorname{conv}\{v\pm\hat e_x/2,\ v\pm\hat e_y/2,\ v\pm\hat e_z/2\}.
\]
The site $v+\hat e_\mu/2$ is shared by $O_v$ and $O_{v+\hat e_\mu}$, giving
exactly two possible owners. Centres are not spin sites. A cubic cell has eight
centres and twelve shared sites on its twelve bonds.

\begin{center}
\includegraphics[width=.94\linewidth]{geometry-origin.pdf}
\end{center}
\vspace{-6pt}
\captionof{figure}{Six tetrahedral edge midpoints form an octahedron. Its eight faces
are four old-face triangles and four vertex-star triangles. At right, blue sites
lie on the bonds between cubic centres (orange).}

Replacing each tetrahedron by its \emph{inscribed} octahedron keeps the six sites
and the specified cubic ownership graph. This differs from exchanging the original
cells of a tetrahedral--octahedral tiling: its original octahedron centres form FCC.

\textbf{Local objects.} Allow any subset of the three opposite-site diameter dimers
$D_x,D_y,D_z$ ($1+3+3+1=8$ possibilities), or exactly one of the eight face trimers
$T_{s_xs_ys_z}$ ($s_\mu=\pm$). Every physical site is covered once globally;
its cubic-bond arrow points toward its owner.

\begin{center}
\includegraphics[width=.94\linewidth]{octahedral-catalogue.pdf}
\end{center}
\vspace{-6pt}
\captionof{figure}{The sixteen allowed local supports. Filled sites are owned by
the displayed octahedron; gray sites belong to neighbours. Crossing diameter
chords meet at a centre, which carries no spin.}

Two complementary face trimers are \textbf{forbidden}. Otherwise all six incoming
sites would have five realizations: three diameters, or one of four complementary
face pairs. With the exclusion, each support specifies its objects uniquely.
For incoming bits $a^{\pm}_{v,\mu}$ the sixteen supports are exactly
\[
 a^+_{v,x}\xor a^-_{v,x}
 =a^+_{v,y}\xor a^-_{v,y}
 =a^+_{v,z}\xor a^-_{v,z}.
\]

\newpage
\sect{4. Three-dimensional model (continued): cage nets and singlets}
Fix a full reference covering $R$ and set
$n_e=a_{v,e}(C)\xor a_{v,e}(R)$. Both ownership bits complement at the other end
of a cubic bond, so $n_e$ is endpoint-independent. Their differences obey
\[
 p_x=p_y=p_z,\qquad p_\mu=n_{+\mu}\xor n_{-\mu}.
\]
This is even incidence in every coordinate plane: an \textbf{X-cube cage net}.
Conversely, reversing reference arrows on any such net preserves one owner per
site and the sixteen-support rule. Unique objects at each support give a covering.
Thus coverings and cage nets are in bijection, for a fixed reference and compatible
boundaries, as in the tetrahedral representation.

For the subsequent quantum extension put \textbf{spin 1} at every site. In the Cartesian basis,
the normalized two- and three-spin singlets are
\[
 d_{\alpha\beta}=\delta_{\alpha\beta}/\sqrt3,
 \qquad t_{\alpha\beta\gamma}=\eps_{\alpha\beta\gamma}/\sqrt6,
 \qquad \ket{\Psi_{3D}}=Z^{-1/2}\sum_{C\ {\rm allowed}}\ket{C}.
\]
Order each trimer by its $x,y,z$ vertices, independent of face signs, with
$\eps_{xyz}=+1$. Dimers are symmetric and trimers antisymmetric under exchange.
Here $\ket{C}$ is a product of normalized object singlets and
$Z=\sum_{C,C'}\langle C'|C\rangle$, including their nonorthogonal overlaps.
An exact virtual-dimension-four tensor network represents this trial state;
no parent Hamiltonian is assumed. Finite-box spin calculations reach $6^3$
centres (648 spins), but \textbf{the 3D bulk correlation length remains unresolved}.

\sect{5. An $xy$ layer viewed from above: superposed dummy-link sites}
\begin{center}
\includegraphics[width=\linewidth]{layer-relation.pdf}
\end{center}
\vspace{-6pt}
\captionof{figure}{Two distinct apices $z^+,z^-$ project onto the same orange point.
Opening this double point displays a dummy link and the decorated square--octagon
layer. Green dashed triangles are selected face \emph{frames}, not two occupied trimers.}

Take centres $v=(m,n,0)$. The four equatorial sites project to
$v\pm\hat e_x/2,v\pm\hat e_y/2$, while both $v\pm\hat e_z/2$ project to $v$:
\textbf{the two dummy-link sites are on top of each other}. Choose a checkerboard
set of centre-lattice plaquettes. At each $v$, the two selected opposite quadrants
group the horizontal ports into adjacent pairs; assign $z^+$ to the upper pair
and $z^-$ to the lower pair. These two complementary face frames, opened for
display, give the upright gray squares and triangle slots of the 2D geometry.

This establishes the proposed geometric connection. The full 3D model still
allows all eight face trimers and opposite-site dimers; the 2D state uses triangle
dimers and spin $1/2$. A reduction must therefore specify the state map and
interlayer amplitudes. Neighbouring centre layers share vertical sites:
$z^+$ of layer $k$ is $z^-$ of layer $k+1$. The cage rule also couples layers: an $xy$ turn needs odd
vertical incidence. Two matching plaquette loops joined by four vertical edges
form an allowed cube cage; a lone planar plaquette does not. These are useful
starting constraints for a foliation construction, which remains to be proved.

{\footnotesize
[4] O.~Benton and R.~Moessner, \emph{Topological Route to New and Unusual Coulomb Spin
Liquids}, Phys.~Rev.~Lett.~127, 107202 (2021),
\href{https://arxiv.org/abs/2103.10817}{arXiv:2103.10817}.\par
[5] M.~Stern, M.~D.~Burke, M.~J.~P.~Gingras, J.~Romh\'anyi and K.~T.~K.~Chung,
\emph{Fracton Spin Liquid and Exotic Frustrated Phases in Ising-like Octochlore
Magnets} (2026), \href{https://arxiv.org/abs/2603.12313}{arXiv:2603.12313},
Secs.~II,V. Their local ground-state catalogue has ten patterns; the covering
model here has sixteen. These references accompany the original geometric notes.}

\end{document}
